% 2009 年广东卷（文科）第 11 题：统计 6 名队员三分球总数的程序框图。
% 来源 PDF：.work/ocr-local/sources/2009广东文.pdf
% SHA-256：26c26718ce8f1cd1128fca4e63e1576ecec72ee27cc46ad538673bb533e61ced
% 原页：第 1 页（A3 三栏）第 2 栏，印刷页 521。
% 绘制依据：开始 → 输入 n, a1, …, a6 → s=0, i=1 → 判断框（原页留空，待填）：
%   是时 s←s+a_i、i←i+1 后回环；否时输出 s、结束。
\begin{tikzpicture}[semithick, every node/.style={inner sep=3pt},
  box/.style={draw, minimum width=2.1cm, minimum height=0.55cm},
  io/.style={draw, trapezium, trapezium left angle=70, trapezium right angle=110,
             minimum width=2.9cm, minimum height=0.55cm},
  dec/.style={draw, diamond, aspect=2, inner sep=1pt,
              minimum width=2.3cm, minimum height=1.0cm}]

  \node[box, rounded corners=3pt] (start) at (0,0) {开始};
  \node[io] (input) at (0,-1.1) {输入 $n,a_1,a_2,\cdots,a_6$};
  \node[box] (init) at (0,-2.2) {$s=0,\ i=1$};
  \node[dec] (test) at (0,-3.5) {};
  \node[io] (out) at (0,-4.9) {输出 $s$};
  \node[box, rounded corners=3pt] (end) at (0,-6.0) {结束};

  \node[box] (assign) at (3.8,-3.5) {$s=s+a_i$};
  \node[box] (inc) at (3.8,-2.5) {$i=i+1$};

  \draw[->,>=stealth] (start) -- (input);
  \draw[->,>=stealth] (input) -- (init);
  \draw[->,>=stealth] (init) -- (test);
  \draw[->,>=stealth] (test) -- node[left] {否} (out);
  \draw[->,>=stealth] (out) -- (end);
  \draw[->,>=stealth] (test.east) -- node[above] {是} (assign.west);
  \draw[->,>=stealth] (assign.north) -- (inc.south);
  \draw[->,>=stealth] (inc.west) -- (0,-2.5);
\end{tikzpicture}
